% =====================================================================
\chapter{観測量ごとの外れと prior の物理}
\label{ch:priors}
% =====================================================================

前章で、利得を決めるのは観測量ごとの外れ $\Delta$ であることを見た。この章では、ハバード模型の現実的な prior について $\Delta$ が観測量ごとにどう振る舞うか、それがなぜかを調べる。データは厳密な参照を持つ 80 系(第\ref{sec:reference}節)の中央のサイトとその隣で、利得は $n_m=100$ での厳密な値である。

\section{全体像}

表\ref{tab:overview}に、$U=12$ の 20 系での利得の中央値をまとめる。

\begin{table}[t]
\centering\small
\caption{$U=12$、厳密な参照を持つ 20 系での利得の中央値($n_m=100$)。二重占有は 4 項の和として測るときの厳密な値。「天井」は prior が完全な場合の値。}
\label{tab:overview}
\begin{tabular}{@{}lccccccc@{}}
\toprule
観測量 & UHF & UHF-sym & PHF & 切断 $\chi{=}4$ & 変分 $\chi{=}8$ & 変分 $\chi{=}8$-sym & 天井 \\
\midrule
オンサイト $ZZ$(電荷) & 465 & 465 & 513 & 151 & 555 & 555 & 555 \\
隣接 $Z_\up Z_\up$(スピン) & 1.2 & 7.0 & 14.2 & 29.2 & 16.8 & 19.8 & 31.6 \\
単一サイト $Z_\up$ & 0.02 & 1.00 & 1.00 & 1.00 & 0.11 & 1.00 & 1 \\
二重占有(4 項の和) & 1.8 & 123 & 127 & 78 & 13.3 & 129 & --- \\
\bottomrule
\end{tabular}
\end{table}

この表から 3 つのことが読み取れる。
\begin{enumerate}
  \item \textbf{観測量によって prior の得意・不得意が逆転する。} UHF は電荷の量(オンサイト $ZZ$)では天井の 8 割に届くが、スピンの量では損をする。切断 MPS($\chi=4$)はその逆である。
  \item \textbf{対称性を回復すると損が消える。} UHF も変分 MPS も、スピン回転で平均すると単一サイト $Z_\up$ と二重占有の損が消える。
  \item \textbf{電荷とスピンの両方で良いのは、変分 MPS を回転平均したものである。}
\end{enumerate}
以下、それぞれの理由を説明する。

\section{電荷の量:強結合で平均場の誤差が消える}
\label{sec:charge}

\subsection{二重占有の誤差は 28\% で止まる}
強結合での真の二重占有率は、式\eqref{eq:d-strong}より 1 次元で
\begin{equation}
  d_{\rm 真}\approx\frac{4t^2}{U^2}\Bigl(\frac14-\ev{\Sv_i\cdot\Sv_j}\Bigr)
\end{equation}
で、隣の結合が一重項を組む確率に比例する($U=12$ で予測 $0.0188$、実測 $0.0180$)。一方、共線 UHF は $\up$ と $\dn$ が独立な Slater 行列式の積なので
\begin{equation}
  d_{\rm UHF}=\ev{n_{i\up}}\ev{n_{i\dn}}=\frac14-m^2
\end{equation}
($m$ は副格子磁化)である。自己無撞着条件を $U\gg t$ で展開すると $m\approx\frac12-\frac{2t^2}{U^2}$ となり
\begin{equation}
  d_{\rm UHF}\approx\frac{2t^2}{U^2}=\frac{4t^2}{U^2}\Bigl(\frac14+\frac14\Bigr).
\end{equation}
括弧の中は古典的な反平行の積状態 $\ket{\up\dn}$ の一重項成分 $\frac14-(-\frac14)=\frac12$、つまり「半分だけ一重項」である。真の値は無限の Heisenberg 鎖の $\ev{\Sv_i\cdot\Sv_j}=\frac14-\ln2$ を使って $\frac14+(\ln2-\frac14)=0.693$ なので、二重占有の相対誤差は
\begin{equation}
  \delta\equiv\frac{d_{\rm UHF}}{d_{\rm 真}}-1\longrightarrow\frac{0.5}{0.693}-1=-0.28
\end{equation}
と、$U$ によらない定数に止まる。

\subsection{オンサイト $ZZ$ に対する誤差は $(t/U)^2$ で消える}
ところが利得を決めるのは、オンサイト $ZZ$ $=4d-1$ に対する相対誤差である:
\begin{equation}
  \varepsilon_{ZZ}=\frac{4\lvert d_{\rm 真}-d_{\rm UHF}\rvert}{\lvert4d_{\rm 真}-1\rvert}\approx4d_{\rm 真}\lvert\delta\rvert\approx3.1\Bigl(\frac tU\Bigr)^2.
\end{equation}
分子は $d\propto(t/U)^2$ とともに 0 へ向かい、分母は 1 に張り付く。\textbf{二重占有そのものの誤差は 28\% のまま残るのに、オンサイト $ZZ$ に対する誤差は $(t/U)^2$ で消える}。これが「平均場 prior は電荷の量では強結合で良くなる」ことの正体である。同時に、天井 $G_{\max}\propto x^2/(1-x^2)$ は $1-x^2\approx8d\propto(t/U)^2$ で上がるので、オンサイト $ZZ$ の利得は強結合で大きくなる。

\subsection{弱結合側}
弱結合では、真の状態は相関で二重占有を抑えるのに、UHF は $\up$ と $\dn$ が無相関なので二重占有を過大評価する。$U\simeq4$ で過大評価と過小評価が入れ替わり、誤差がいったん小さくなる。

\subsection{電荷の量が境界条件に鈍い理由}
半充填の二部格子では、粒子正孔対称性によりサイトごとの電子数が厳密に $\ev{n_i}=1$ になる(付録\ref{app:ph})。開放端の影響は電荷の密度には現れず、結合の強弱の交替として現れる。磁化を持つ UHF も、粒子正孔変換とスピン回転を組み合わせた変換の対称性を保つので $\ev{n_i}=1$ を満たす。電荷の量で UHF が境界条件や系の大きさによらず良い prior であるのはこのためである。二部格子でない $4\times3$ トーラスでは UHF の電荷が 6.5\% 偏っていた。

\section{スピンの量:平均場が損をする 2 つの理由}
\label{sec:spin}

\subsection{理由 1:スピンの向きを選ぶ}
有限系の真の基底状態はスピン回転で不変で、$\ev{S^z_i}=0$ である。共線 UHF は副格子ごとに $\ev{S^z_i}=\pm m$ を持つので、単一サイトの $Z_{i\up}=1-n_i-2S^z_i$ を大きく外す。真の値が 0 なので損得の境目を必ず越え、$U=12$ では 20 系すべてで $G=0.02$ 程度の大損になる。隣接 $Z_\up Z_\up$ でも、スピンの向きを選ぶことで $z$ 方向の相関を過大評価し、80 系中 17 系で損をする。

この原因は、スピン回転で平均した UHF-sym(第\ref{sec:rotavg}節)で取り除ける。単一サイト $Z_\up$ は $1-\ev{n_i}$ になって外れが消え、隣接 $Z_\up Z_\up$ で損をする系は 17 から 0 になる。

\subsection{理由 2:量子的な一重項の相関が足りない}
回転平均は $\ev{\Sv_i\cdot\Sv_j}$ を変えない。UHF の $\ev{\Sv_i\cdot\Sv_j}$ は古典的な Néel 配置の値 $-1/4$ で、真の値($1$ 次元の無限鎖で $\frac14-\ln2=-0.443$)の量子的な一重項の相関を持たない。隣接 $\ev{S^z_iS^z_j}$ について、真の値は $\ev{\Sv_i\cdot\Sv_j}/3$、共線 UHF は $-m^2\to-1/4$ なので
\begin{equation}
  \frac yx\longrightarrow\frac{3/4}{\lvert\ev{\Sv_i\cdot\Sv_j}\rvert},\qquad
  \varepsilon\longrightarrow\Bigl\lvert1-\frac{3/4}{\lvert\ev{\Sv_i\cdot\Sv_j}\rvert}\Bigr\rvert
\end{equation}
となる。1 次元では $\varepsilon\to0.69$ で、実測($L=12$ の周期境界)の $0.690$($U=8$)、$0.684$($U=12$)と一致する。損得の境目 $y/x=2$ は $\lvert\ev{\Sv_i\cdot\Sv_j}\rvert=3/8$ に当たり、1 次元(0.443)では得、2 次元の正方格子(約 0.335)では損になる。

この原因は回転平均では直らない。UHF-sym の隣接 $Z_\up Z_\up$ の利得は $U=12$ で 7.0 と、天井 31.6 に遠く届かない。

\section{スピン射影 HF:効果は $1/L$ で消える}
\label{sec:phf}

量子的な相関を足す方法として、UHF を全スピン 0 に射影した PHF を試した。表\ref{tab:overview}のとおり、PHF の隣接 $Z_\up Z_\up$ は 14.2 と UHF-sym の 2 倍だが、天井には届かない。しかも\textbf{この改善は系の大きさとともに消える}。

射影の効果は回転子の描像で説明できる。UHF の 2 つの副格子のスピンはそれぞれ大きさ $\sim mL/2$ の大きなスピンにまとまっている。射影はそれらを全スピン 0 に組み直すが、1 つの結合に割り当てられる量子的な相関は $1/L$ に比例する。実際、射影による隣の $\ev{\Sv_i\cdot\Sv_j}$ の変化は
\begin{equation}
  L\,\Delta\!\ev{\Sv_i\cdot\Sv_j}=-2m
\end{equation}
で予測でき($U=12$、$L=64$ で実測 $-0.963$ 対予測 $-0.972$)、エネルギーの変化も $E_{\rm PHF}-E_{\rm UHF}=-2mJ$ と $L$ によらない。$L=64$ では、UHF と真の値の差のうち射影で埋まるのは 11\% だけである。数十〜数百量子ビットの系で PHF を使う利点はない。

\section{行列積状態の prior:作り方で性質が逆になる}
\label{sec:mps}

低い結合次元 $\chi$ の MPS は、結合ごとに量子的な相関を持てる。ところが、切断 MPS と変分 MPS では性質が逆になった(表\ref{tab:overview})。
\begin{itemize}
  \item \textbf{切断 MPS($\chi=4$)}:隣接 $Z_\up Z_\up$ の利得は 29.2 と天井の 9 割に届くが、二重占有は 78、オンサイト $ZZ$ は 151 にとどまる。
  \item \textbf{変分 MPS($\chi=4$)}:オンサイト $ZZ$ は 553 と天井に届き、相関エネルギーの 78〜88\% を $L$ によらず取り戻すが、中央サイトの $\lvert\ev{S^z}\rvert$ が約 0.30 あり、単一サイト $Z_\up$ では 20 系すべてで損をする。
\end{itemize}
中央の結合で量子数ごとに Schmidt 分解すると、この違いは少ない状態の使い方の違いだと分かる。厳密な状態($L=64$、$U=12$)では電荷がゆらいだ成分の重みが 1.7\% ある。切断 MPS は重みの大きいスピンの成分だけを上下対称に残し、この電荷の成分を捨てる。変分 MPS は電荷の成分を残してその重み(1.8\%)を正しく再現する代わりに、スピンの成分を一方の向きだけにする。

\section{変分 MPS の回転平均:電荷とスピンの両方で良い prior}
\label{sec:varsym}

変分 MPS の弱点はスピンの向きを選ぶことだけなので、UHF-sym と同じ回転平均(式\eqref{eq:rotavg})で直せるはずである。clara に保存した変分 MPS(80 系 × 4 つの $U$ × $\chi=2,4,8,16$ の 320 個)について、回転平均に要る $\ev{n_i}$、$\ev{n_in_j}$、$\ev{\Sv_i\cdot\Sv_j}$ を求めた。

結果を表\ref{tab:varsym}に示す。
\begin{table}[t]
\centering\small
\caption{変分 MPS とその回転平均(-sym)の利得の中央値($n_m=100$、各 $U$ の 20 系)。括弧は損をする系の数。}
\label{tab:varsym}
\begin{tabular}{@{}llcccc@{}}
\toprule
$U$ & 観測量 & UHF-sym & 変分 $\chi{=}8$ & 変分 $\chi{=}8$-sym & 変分 $\chi{=}16$-sym \\
\midrule
4  & 隣接 $Z_\up Z_\up$ & 11.2 & 9.71 (1) & 11.7 & 18.9 \\
4  & 二重占有 & 28.0 & 8.67 (3) & 28.0 & 28.1 \\
12 & 隣接 $Z_\up Z_\up$ & 6.95 & 16.8 (1) & 19.8 & 26.7 \\
12 & 二重占有 & 123 & 13.3 & 129 & 129 \\
\bottomrule
\end{tabular}
\end{table}

\begin{itemize}
  \item $\chi\ge8$ の回転平均は、80 系 × 4 観測量(オンサイト $ZZ$、隣接 $Z_\up Z_\up$、単一サイト $Z_\up$、二重占有)のどれでも損をしなかった(単一サイト $Z_\up$ の 2 系は UHF-sym でも損をする観測量の限界)。
  \item 隣接 $Z_\up Z_\up$ では $U\ge8$ で UHF-sym の 2〜3 倍得をする。回転平均した prior の質は $\ev{\Sv_i\cdot\Sv_j}$ の正確さで決まり、その相対誤差は $U=12$ で UHF 0.35 に対して変分 $\chi=8$ で 0.11、$\chi=16$ で 0.05 だった。
  \item 必要なのは変分 MPS を 1 回作ることと期待値の計算だけで、厳密な状態は要らない。厳密な状態を使わない prior の中では、本研究で試した中で最も良い。
\end{itemize}

\section{和の観測量:項ごとの外れ}
\label{sec:sumobs}

二重占有を 4 項の和\eqref{eq:docc-pauli}として測ると、共線 UHF の利得は $U=12$ で 1.8 しかない(表\ref{tab:overview})。二重占有の合計の外れは小さい(相対 28\%、かつ $d$ 自体が小さい)のに、である。命題\ref{prop:sumvar}によれば、分散には項ごとの外れ $\Delta_k\Delta_l$ が入る。共線 UHF の $\ev{Z_{i\up}}=\pm2m$ は真の値 0 から大きく外れ、$Z_{i\up}$ と $Z_{i\dn}$ の外れは符号が逆で合計では打ち消し合うが、分散には打ち消されずに残る。回転平均すると単一の $Z$ の外れが消え、利得は 123 になる。

\textbf{和の観測量では、合計の値が正しいだけでは足りない。各項の値が正しい必要がある。}逆に、参照状態の側が対称性を破っているとき(2 次元の DMRG の状態など)は、回転平均した prior の方が項ごとの外れが大きくなり、利得が下がることがある(第\ref{ch:2d}章)。

\section{まとめ}
\begin{itemize}
  \item 平均場 prior は、電荷の量(オンサイト $ZZ$)では強結合で誤差が $(t/U)^2$ で消え、天井近くまで得をする。二重占有そのものの誤差は 28\% で止まるが、オンサイト $ZZ$ では分母が 1 に張り付くので効かない。
  \item スピンの量で平均場が損をする原因は 2 つある。スピンの向きを選ぶこと(回転平均で直る)と、量子的な一重項の相関が足りないこと(直らない。2 次元では損になる)。
  \item スピン射影は量子的な相関を足すが、その効果は $1/L$ で消える。
  \item 変分 MPS を回転平均すると、電荷とスピンの両方で損をしない prior になる。
  \item 和の観測量では、項ごとの外れが利得を決める。
\end{itemize}
